\documentclass[acmsmall,nonacm,screen]{acmart}

\usepackage{mathpartir}
\usepackage{booktabs}

\citestyle{acmauthoryear}
\settopmatter{printacmref=false}

\newcommand{\eqd}{\doteq}
\newcommand{\kw}[1]{\mathsf{#1}}
\newcommand{\nm}[1]{\mathit{#1}}
\newcommand{\chk}[4]{#1 \vdash #2 : #3 \eqd #4}
\newcommand{\wf}[1]{\vdash #1}
\newcommand{\consts}{\mathrm{consts}}
\newcommand{\axs}{\mathrm{ax}}
\newcommand{\audit}{\mathrm{audit}}
\newcommand{\restr}[2]{#1{\restriction}_{#2}}
\newcommand{\lean}[1]{\texttt{#1}}

\begin{document}

\title{What Must a Human Read? Audit Surfaces for Machine-Checked Libraries}

\author{Claude-Fable-5.2}
\authornote{An AI model. It wrote this paper, the proofs, and the artifact in one session, at the
request of a human who asked for a first paper for this repository. At the time of writing no
human has reviewed them.}

\begin{abstract}
When an agent proves a theorem and a kernel checks the proof, a human who wants to rely on the
theorem still has to read something. We make precise what. For a library of definitions, axioms,
and theorems, the audit surface of a theorem consists of the definitions that its statement
unfolds to and the axioms that its proof uses, with the definitions those axioms unfold to. We
prove that the statement is derivable from the audit surface alone, and that one pass over the
library computes the surface. Three consequences follow for a library that many agents extend. A
theorem holds in every library that contains its audit surface, independently checked extensions
merge without rechecking, and a definition that its clients never unfold can be replaced
together with the proofs of its laws. All these results are mechanized in Lean~4. We also measure audit
surfaces in Lean's own library. For 2,000 theorems of \lean{Init} and \lean{Std}, the median
audit surface has 1,436 source lines, against 6,241 lines for everything the proof depends on.
\end{abstract}

\keywords{proof assistants, trusted computing base, definitions, axioms, proof libraries}

\authorsaddresses{}

\maketitle

%% 1 ---------------------------------------------------------------------------------------
\section{Introduction}
\label{sec:intro}

Proof assistants divide trust. A small kernel checks every proof, so nobody has to trust the
program or the person that wrote it. The division goes back to LCF~\citep{lcf1979} and is known
as the de Bruijn criterion~\citep{barendregt2005,geuvers2009}. It matters more now that language
models write formal proofs~\citep{jiang2022}: a kernel accepts the proof of an agent on the same
terms as the proof of a person.

The kernel settles whether a proof is correct. It leaves open what was proved. A human who wants
to rely on a theorem must still read, and the usual advice is to read the statement and to skip
the proof~\citep{pollack1998}. The advice leaves out two things. First, a statement mentions
defined constants, and its meaning lies in their definitions, which the same agent may have
written. Second, a proof may use axioms, and an axiom is an assumption of the theorem even when
the statement never mentions it. The safe alternative is to read everything the proof depends on,
which gives up the benefit of the kernel.

This paper determines the set in between. We call it the \emph{audit surface} of a theorem: the
definitions that the statement unfolds to, the axioms that the proof uses, and the definitions
that those axioms unfold to. Lemmas, proofs, and definitions that occur only in proofs stay
outside.

Our main contributions are as follows:
\begin{itemize}
\item A library calculus with definitions, axioms, and theorems that carry proof terms, and a
  kernel of about forty lines (Section~\ref{sec:calculus}).
\item The audit surface, with the theorem that the statement of a checked proof is derivable from
  its audit surface alone, and a one-pass computation of the surface as a least closed set
  (Section~\ref{sec:audit}).
\item Three consequences for a shared library: a theorem holds wherever its audit surface is
  present, extensions with fresh names merge without rechecking, and a definition that its
  clients never unfold can be replaced together with the proofs of its laws
  (Section~\ref{sec:structural}).
\item A measurement of audit surfaces for 2,000 theorems of Lean's library
  (Section~\ref{sec:measure}).
\end{itemize}

\paragraph{Status of the claims.}
Every definition, lemma, proposition, theorem, and corollary of
Sections~\ref{sec:calculus}--\ref{sec:structural} is mechanized in Lean~4~\citep{moura2021}. The
development has about 1,500 lines and uses no axioms beyond the three standard axioms of Lean. The
numbers of Section~\ref{sec:measure} are printed by a Lean program. The command \lean{make} in
the artifact checks the proofs and reproduces the numbers with Lean 4.36.0, in about two minutes. The examples of
Section~\ref{sec:necessity} and the discussion in Section~\ref{sec:limits} are argued informally.

%% 2 ---------------------------------------------------------------------------------------
\section{A Library Calculus}
\label{sec:calculus}

The calculus has untyped $\lambda$-terms with constants, proof terms for equations between
them, and libraries that list definitions, axioms, and theorems.
\[
\begin{array}{llcl}
\text{Terms} & s, t, u & ::= & x \mid c \mid s\;t \mid \lambda x.\,s \\
\text{Proofs} & p, q & ::= & \kw{refl}\;t \mid \kw{sym}\;p \mid \kw{trans}\;p\;q \mid
  \kw{app}\;p\;q \mid \kw{lam}\;p \\
 & & \mid & \kw{beta}\;s\;t \mid \kw{delta}\;c \mid \kw{ax}\;a \mid \kw{lem}\;h \\
\text{Entries} & e & ::= & \kw{def}\;c := u \mid \kw{axiom}\;a : s \eqd t \mid
  \kw{theorem}\;h : s \eqd t := p \\
\text{Libraries} & L & ::= & \cdot \mid L, e
\end{array}
\]
Constants $c$, axiom names $a$, and theorem names $h$ come from one set of names. We write
$\consts(t)$ for the constants of $t$. A constant needs no declaration. A constant without a
definition is a parameter, about which the library says nothing except through axioms.

\begin{figure}
\begin{mathpar}
\inferrule{ }{\chk{L}{\kw{refl}\;t}{t}{t}}
\and
\inferrule{\chk{L}{p}{s}{t}}{\chk{L}{\kw{sym}\;p}{t}{s}}
\and
\inferrule{\chk{L}{p}{s}{u} \\ \chk{L}{q}{u}{t}}{\chk{L}{\kw{trans}\;p\;q}{s}{t}}
\and
\inferrule{\chk{L}{p}{s}{s'} \\ \chk{L}{q}{t}{t'}}{\chk{L}{\kw{app}\;p\;q}{s\;t}{s'\;t'}}
\and
\inferrule{\chk{L}{p}{s}{s'}}{\chk{L}{\kw{lam}\;p}{\lambda x.\,s}{\lambda x.\,s'}}
\and
\inferrule{ }{\chk{L}{\kw{beta}\;s\;t}{(\lambda x.\,s)\;t}{s[t/x]}}
\and
\inferrule{(\kw{def}\;c := u) \in L}{\chk{L}{\kw{delta}\;c}{c}{u}}
\and
\inferrule{(\kw{axiom}\;a : s \eqd t) \in L}{\chk{L}{\kw{ax}\;a}{s}{t}}
\and
\inferrule{(\kw{theorem}\;h : s \eqd t := q) \in L}{\chk{L}{\kw{lem}\;h}{s}{t}}
\end{mathpar}
\Description{Nine inference rules that define the checking judgment.}
\caption{Checking. The proof $p$ establishes $s \eqd t$ in the library $L$.}
\label{fig:chk}
\end{figure}

\begin{figure}
\begin{mathpar}
\inferrule{ }{\wf{\cdot}}
\and
\inferrule{\wf{L} \\ c \text{ fresh for } L \\ u \text{ closed} \\ c \notin \consts(u)}
  {\wf{L,\, \kw{def}\;c := u}}
\and
\inferrule{\wf{L} \\ a \text{ fresh for } L}{\wf{L,\, \kw{axiom}\;a : s \eqd t}}
\and
\inferrule{\wf{L} \\ h \text{ fresh for } L \\ \chk{L}{p}{s}{t}}
  {\wf{L,\, \kw{theorem}\;h : s \eqd t := p}}
\end{mathpar}
\Description{Four inference rules that define the well-formed libraries.}
\caption{Well-formed libraries.}
\label{fig:wf}
\end{figure}

\paragraph{Checking.}
Figure~\ref{fig:chk} defines when a proof $p$ establishes an equation in a library, written
$\chk{L}{p}{s}{t}$. The first five rules make $\eqd$ a congruence. The rule for $\kw{beta}$
contracts a $\beta$-redex, and the rule for $\kw{delta}$ unfolds a definition. The last two rules
use an axiom and a theorem of the library. A theorem is used through its statement only.

\paragraph{Libraries.}
Figure~\ref{fig:wf} defines the well-formed libraries. A name is \emph{fresh} for $L$ when no
entry of $L$ has this name and the name occurs in no body and no statement of $L$. The body of a
definition is closed and does not mention the constant it defines. With freshness, this rules out
recursive definitions. A theorem is checked against the entries before it.

The rules of Figure~\ref{fig:chk} consult the library by membership only, which gives the first
structural property.

\begin{lemma}[Locality]
\label{lem:local}
Let $\chk{L}{p}{s}{t}$. Let $L'$ contain every definition of $L$ that $p$ unfolds, every axiom
of $L$ that $p$ uses, and, for every theorem of $L$ that $p$ uses, a theorem with the same name
and statement. Then $\chk{L'}{p}{s}{t}$. In particular, a proof that checks in a library checks
in every library that contains its entries.
\end{lemma}

\begin{figure}
\[
\begin{array}{ll}
\kw{def} & \nm{flip} := \lambda f.\,\lambda x.\,\lambda y.\,f\;y\;x \\
\kw{axiom} & \nm{add\_comm} : \nm{flip}\;\nm{add} \eqd \nm{add} \\
\kw{def} & \nm{double} := \lambda x.\,\nm{add}\;x\;x \\
\kw{def} & \nm{comp} := \lambda f.\,\lambda g.\,\lambda x.\,f\;(g\;x) \\
\kw{theorem} & \nm{comp\_double} : \nm{comp}\;\nm{double}\;\nm{double} \eqd
  \lambda x.\,\nm{add}\;(\nm{add}\;x\;x)\;(\nm{add}\;x\;x) \\
\kw{def} & \nm{quad} := \lambda x.\,\nm{double}\;(\nm{double}\;x) \\
\kw{theorem} & \nm{quad\_comp} : \nm{quad} \eqd \nm{comp}\;\nm{double}\;\nm{double} \\
\kw{theorem} & \nm{quad\_eq} : \nm{quad} \eqd
  \lambda x.\,\nm{add}\;(\nm{add}\;x\;x)\;(\nm{add}\;x\;x) \\
\kw{theorem} & \nm{add\_flip} : \nm{add} \eqd \lambda x.\,\lambda y.\,\nm{add}\;y\;x \\
\kw{theorem} & \nm{double\_sum} : \lambda x.\,\lambda y.\,\nm{double}\;(\nm{add}\;x\;y) \eqd
  \lambda x.\,\lambda y.\,\nm{add}\;(\nm{add}\;y\;x)\;(\nm{add}\;x\;y)
\end{array}
\]
\Description{A library with four definitions, one axiom, and five theorems.}
\caption{A library, oldest entry first. The proofs are described in the text.}
\label{fig:example}
\end{figure}

\paragraph{Example.}
Figure~\ref{fig:example} shows a library with four definitions, one axiom, and five theorems.
The constant $\nm{add}$ has no definition. The axiom $\nm{add\_comm}$ states that it is
commutative, in terms of the defined constant $\nm{flip}$. The theorems $\nm{comp\_double}$ and
$\nm{quad\_comp}$ are proved by normalizing both sides. The proof of $\nm{quad\_eq}$ is
$\kw{trans}\;(\kw{lem}\;\nm{quad\_comp})\;(\kw{lem}\;\nm{comp\_double})$. The proof of
$\nm{add\_flip}$ is $\kw{trans}\;(\kw{sym}\;(\kw{ax}\;\nm{add\_comm}))\;q$, where $q$ normalizes
$\nm{flip}\;\nm{add}$. The proof of $\nm{double\_sum}$ normalizes the left side, rewrites one
occurrence of $\nm{add}$ with $\kw{lem}\;\nm{add\_flip}$, and normalizes again.

\paragraph{A kernel and an untrusted prover.}
The mechanization represents terms with de Bruijn indices, lists a library newest entry first,
and represents a set of names by a Boolean predicate. It defines a function
$\mathrm{infer}_L(p)$ that computes the equation a proof establishes, if there is one. This
function is the kernel. A second function decides the premises of Figure~\ref{fig:wf} with the
kernel.

\begin{proposition}[Kernel soundness]
\label{prop:kernel}
If $\mathrm{infer}_L(p) = (s, t)$, then $\chk{L}{p}{s}{t}$. If the decision procedure accepts
$L$, then $\wf{L}$.
\end{proposition}

The mechanization also contains a prover that normalizes terms and returns proof terms. Nothing is
proved about the prover. Its proofs enter a library only when the kernel accepts them, and this
is how the library of Figure~\ref{fig:example} is checked. The division of labor is the one
between an agent and a kernel.

%% 3 ---------------------------------------------------------------------------------------
\section{Audit Surfaces}
\label{sec:audit}

\subsection{Definition}

The audit surface depends on the axioms that a proof uses, directly or through theorems. For a
proof $p$ and a library $L$, replace every $\kw{lem}\;h$ in $p$ by the proof of $h$ in $L$, for
each theorem $h$ of $L$ from the newest to the oldest. We write $\axs_L(p)$ for the set of
names $a$ such that $\kw{ax}\;a$ occurs in the result.

\begin{definition}[Closed set]
\label{def:closed}
A set $S$ of names is \emph{closed} for $L$ when $\consts(u) \subseteq S$ for every definition
$\kw{def}\;c := u$ of $L$ with $c \in S$, and $\consts(s) \cup \consts(t) \subseteq S$ for
every axiom $\kw{axiom}\;a : s \eqd t$ of $L$ with $a \in S$.
\end{definition}

\begin{definition}[Audit set and audit surface]
\label{def:audit}
Let $\chk{L}{p}{s}{t}$. The \emph{audit set} of $p$ is the least set that is closed for $L$ and
contains $\consts(s)$, $\consts(t)$, and $\axs_L(p)$. For a set $S$ of names, the restriction
$\restr{L}{S}$ consists of the definitions and the axioms of $L$ whose names are in $S$, in the
order of $L$. The \emph{audit surface} of $p$ is $\restr{L}{S}$ for the audit set $S$ of $p$.
\end{definition}

The audit set is computed without iteration.

\begin{proposition}[One pass]
\label{prop:pass}
Let $\wf{L}$ and let $A$ be a set of names. Start with $A$ and go through the entries of $L$
from the newest to the oldest. Whenever the name of a definition or of an axiom is in the current
set, add the constants of its body or of its statement. The result is the least set that contains
$A$ and is closed for $L$.
\end{proposition}

One pass suffices because of freshness: the name of an entry occurs in no older entry, so nothing
adds the name after the pass has moved beyond the entry.

In Figure~\ref{fig:example}, the statement of $\nm{quad\_eq}$ mentions $\nm{quad}$ and
$\nm{add}$, and its proof uses no axiom. The pass adds $\nm{double}$ at the definition of
$\nm{quad}$. The audit surface consists of the definitions of $\nm{quad}$ and $\nm{double}$. It
contains neither $\nm{comp}$ nor the two lemmas, although the proof uses both lemmas and their
statements mention $\nm{comp}$. The statement of $\nm{double\_sum}$ mentions $\nm{double}$ and
$\nm{add}$. Its proof uses $\nm{add\_flip}$, whose proof uses the axiom $\nm{add\_comm}$. The
audit surface consists of $\nm{double}$, $\nm{add\_comm}$, and $\nm{flip}$.

\subsection{Soundness}

\begin{theorem}[Audit soundness]
\label{thm:audit}
Let $\wf{L}$ and $\chk{L}{p}{s}{t}$. Let $S$ be closed for $L$ and contain $\consts(s)$,
$\consts(t)$, and $\axs_L(p)$. Then $\restr{L}{S}$ is well formed and contains no theorem, and
there is a proof $p'$ such that $\chk{\restr{L}{S}}{p'}{s}{t}$. In particular, the statement of
a checked proof is derivable from its audit surface.
\end{theorem}

\begin{proof}[Proof outline]
The proof $p'$ is computed from $p$ in two passes over $L$, each from the newest entry to the
oldest. The first pass removes the theorems. A theorem $h$ of $L$ was checked against the entries
before it, so its proof $q$ checks in $L$ by Lemma~\ref{lem:local}. Replacing $\kw{lem}\;h$ by
$q$ gives a proof of the same equation, because names are unique. After the pass the proof checks
against the definitions and the axioms of $L$.

The second pass removes the definitions and the axioms outside $S$. For a definition
$\kw{def}\;c := u$ with $c \notin S$, substitute $u$ for $c$ in every term of the proof and
replace $\kw{delta}\;c$ by $\kw{refl}\;u$. The result is a proof. Substituting a closed term
commutes with $\beta$-contraction. A step $\kw{delta}\;d$ for another definition
$\kw{def}\;d := u_d$ is unchanged because $c$ does not occur in $u_d$. If $d$ is older than
$c$, this is the freshness of $c$. If $d$ is newer, then $d \in S$, since the pass has already
removed the newer definitions outside $S$, and $S$ is closed. The same argument applies to the
axioms that the proof uses, which are in $S$. The statement is unchanged because its constants
are in $S$. An axiom outside $S$ does not occur in the proof and is dropped by
Lemma~\ref{lem:local}.
\end{proof}

\subsection{Each Part Is Needed}
\label{sec:necessity}

The three parts of the audit surface are needed in the following sense. The arguments use that
without axioms the equations derivable in a library are the $\beta\delta$-conversions, and that
terms with distinct normal forms are not convertible~\citep{church1936}. The arguments are not
mechanized.

\emph{Definitions in the statement.} Remove the definition of $\nm{double}$ from the audit
surface of $\nm{quad\_eq}$. In the library that only defines $\nm{quad}$, the two sides of
$\nm{quad\_eq}$ have the distinct normal forms $\lambda x.\,\nm{double}\;(\nm{double}\;x)$ and
$\lambda x.\,\nm{add}\;(\nm{add}\;x\;x)\;(\nm{add}\;x\;x)$.

\emph{Axioms that the statement does not mention.} Remove $\nm{add\_comm}$ from the audit surface
of $\nm{double\_sum}$. The two sides then have normal forms that differ in one argument, namely
$\nm{add}\;x\;y$ and $\nm{add}\;y\;x$. The kernel rejects the proof that only normalizes, and the
mechanization records this.

\emph{Definitions in axioms.} Change the body of $\nm{flip}$ to $\lambda f.\,f$. The axiom
$\nm{add\_comm}$ keeps its name and its text, and it now states a $\beta\delta$-conversion. The
statement of $\nm{add\_flip}$ is no longer derivable. A reader who skips $\nm{flip}$ cannot tell
the two libraries apart.

\subsection{Subjects and Vocabulary}
\label{sec:subjects}

A reader may leave some bodies of the audit surface unread. The theorem then says something
about those constants as artifacts that are identified by name. For example, a reader of
$\nm{double\_sum}$ who skips the body of $\nm{double}$ learns that the function called
$\nm{double}$ in this library, whatever it computes, satisfies the stated equation. This reading
is the one a client of an implementation wants, in the way a client of an abstract type relies on
its operations without their representation~\citep{mitchell1988}. We call such constants the
\emph{subjects} of the theorem and the remaining ones its \emph{vocabulary}. The reading is sound
by Theorem~\ref{thm:audit}, since the bodies are in the audit surface whether they are read or
not. The choice of subjects belongs to the reader. The \emph{signature surface}, in which no body
is read, bounds the reading effort from below, and Section~\ref{sec:measure} measures both.

%% 4 ---------------------------------------------------------------------------------------
\section{Reading, Merging, Replacing}
\label{sec:structural}

A library that many agents extend supports three operations. A reviewer reads a theorem, two
contributions are merged, and a part is replaced by a better one. Each should cost no rechecking
of what is unaffected.

\begin{corollary}[Portability]
\label{cor:portable}
Let $\wf{L}$ and $\chk{L}{p}{s}{t}$, and let $p'$ be the proof of Theorem~\ref{thm:audit} for
the audit set of $p$. Then $\chk{L'}{p'}{s}{t}$ for every library $L'$ that contains the audit
surface of $p$.
\end{corollary}

A reviewer who has read an audit surface has reviewed the theorem for every library in which the
surface is present. The library may replace all proofs and all definitions outside the surface.

\begin{theorem}[Merging]
\label{thm:merge}
Let $\wf{L_0, L_1}$ and $\wf{L_0, L_2}$, and let every name that $L_1$ introduces be fresh for
$L_2$. Then $\wf{L_0, L_2, L_1}$.
\end{theorem}

The proof uses Lemma~\ref{lem:local} and checks no proof again. The condition in the other
direction is not needed for well-formedness. It matters for meaning: a constant that $L_1$ uses
as a parameter and $L_2$ defines is captured by the merge. Systems in which every constant is
declared exclude this case. By Corollary~\ref{cor:portable}, every statement proved in $L_0, L_1$
or in $L_0, L_2$ stays derivable in the merged library from the audit surface it had before.

Two libraries have the \emph{same interface} when they have the same entries in the same order,
except that the proofs of corresponding theorems may differ.

\begin{theorem}[Replacement]
\label{thm:replace}
Let $\wf{L_0,\, \kw{def}\;c := u,\, M,\, L_1}$ and $\wf{L_0,\, \kw{def}\;c := u',\, M'}$, where
$M$ and $M'$ have the same interface. Suppose that no proof of a theorem of $L_1$ contains
$\kw{delta}\;c$, and that no name that $L_1$ introduces occurs in $u'$. Then
$\wf{L_0,\, \kw{def}\;c := u',\, M',\, L_1}$.
\end{theorem}

The theorem covers the repair of an implementation. The entries $M$ state the laws of $c$, with
proofs that unfold it. The entries $L_1$ are clients that use the laws and never unfold $c$. A
new body $u'$ with new proofs of the same laws leaves every client checked. A variant without the
definition, which is also mechanized, says that the proofs of theorems can always be replaced.

%% 5 ---------------------------------------------------------------------------------------
\section{Audit Surfaces in Lean's Library}
\label{sec:measure}

We measure how much a reader saves by reading audit surfaces in a real library.

\paragraph{Method.}
The kernel of Lean~4 has more kinds of declarations than our calculus, and we extend the
definitions as follows. A constant that is read as part of a statement brings in the constants of
its type. A definition also brings in the constants of its body, an inductive type brings in its
constructors, and a constructor or a recursor brings in its inductive type. A theorem or an opaque
constant that occurs in a statement contributes its type only. The axioms of a proof are the
axioms among the constants it depends on, as \lean{\#print axioms} reports them. For a theorem we
compute three sets of constants:
\begin{itemize}
\item the \emph{signature surface}, reached from the statement and the axioms without the body of
  any definition;
\item the \emph{audit surface}, reached with the bodies of definitions;
\item the \emph{dependency closure}, reached with the bodies of definitions and opaque constants
  and with the proofs of theorems, which is what the kernel checked.
\end{itemize}
We report the number of constants and the number of distinct source lines of the declarations
that introduce them. Lean records a source range for each declaration and attributes recursors
and similar generated constants to the declaration of their type. A generated constant without a
range adds no lines.

The population consists of the 46,992 theorems of the modules under \lean{Init} and \lean{Std}
in Lean 4.36.0 that have a source range of their own and a name that is not internal. The sample
takes every 23rd theorem in the order of names, which gives 2,000 theorems.

\begin{table}
\caption{Sizes for 2,000 theorems of \lean{Init} and \lean{Std}.}
\label{tab:sizes}
\begin{tabular}{llrrrrrr}
\toprule
 & & p25 & median & p75 & p90 & max & mean \\
\midrule
Constants & signature surface & 33 & 64 & 153 & 184 & 247 & 87 \\
 & audit surface & 123 & 251 & 659 & 764 & 1,804 & 367 \\
 & dependency closure & 353 & 1,819 & 3,578 & 3,911 & 7,102 & 1,948 \\
\midrule
Source lines & signature surface & 222 & 462 & 1,020 & 1,209 & 1,688 & 595 \\
 & audit surface & 637 & 1,436 & 3,107 & 3,340 & 6,758 & 1,749 \\
 & dependency closure & 1,255 & 6,241 & 11,045 & 11,641 & 23,729 & 6,181 \\
\bottomrule
\end{tabular}
\end{table}

\begin{table}
\caption{Constants and source lines for selected theorems.}
\label{tab:examples}
\begin{tabular}{lrrrrrr}
\toprule
 & \multicolumn{2}{c}{Signature} & \multicolumn{2}{c}{Audit} & \multicolumn{2}{c}{Closure} \\
Theorem & const. & lines & const. & lines & const. & lines \\
\midrule
\lean{Nat.add\_comm} & 13 & 88 & 26 & 120 & 47 & 224 \\
\lean{List.reverse\_reverse} & 9 & 88 & 21 & 112 & 37 & 200 \\
\lean{Array.size\_push} & 27 & 209 & 54 & 286 & 68 & 352 \\
\lean{Nat.gcd\_comm} & 12 & 153 & 138 & 753 & 487 & 1,755 \\
\lean{String.length\_append} & 71 & 468 & 445 & 2,308 & 3,490 & 10,731 \\
\lean{Std.HashMap.size\_insert} & 162 & 1,088 & 621 & 3,177 & 3,371 & 11,147 \\
\lean{Std.TreeMap.contains\_insert} & 180 & 1,174 & 726 & 2,954 & 3,743 & 10,587 \\
\bottomrule
\end{tabular}
\end{table}

\paragraph{Results.}
Table~\ref{tab:sizes} gives the distribution of the three sizes, and Table~\ref{tab:examples}
gives individual theorems. We observe the following.

\emph{The audit surface is a fraction of what the kernel checked.} The dependency closure of a
theorem is larger than its audit surface by a median factor of 4.9 in constants and 3.3 in source
lines. For 1,509 of the 2,000 theorems, the audit surface has at most half the lines of the
closure.

\emph{The audit surface is large.} Its median is 1,436 lines, and it is at most 500 lines for 432
of the 2,000 theorems. A statement is a line or two, so the advice to read the statement
understates the reading by about three orders of magnitude.

\emph{Bodies account for a factor of 2.8.} The audit surface has a median of 2.8 times the lines
of the signature surface. Theorems about data structures show the difference most. The statement
of \lean{Std.HashMap.size\_insert} mentions the insertion function, so its audit surface holds
the implementation of hash maps, with 3,177 lines against 1,088 for the signatures. This is the
case of Section~\ref{sec:subjects}: a client treats the insertion function as a subject.

\emph{Surfaces overlap.} The union of the 2,000 audit surfaces has 8,328 constants and 32,768
lines, against 31,319 constants and 101,918 lines for the union of the closures. Reading the
union once audits all 2,000 theorems, which is about 16 lines per theorem.

\emph{Axioms are few.} The sample uses only the three standard axioms of Lean, which are
\lean{propext}, \lean{Classical.choice}, and \lean{Quot.sound}. No axiom is used by 205 of the
theorems.

\paragraph{Threats to validity.}
The measurement applies the definitions to Lean's kernel terms, and we have not proved
Theorem~\ref{thm:audit} for Lean's type theory. Source lines cover whole declarations, so they
overstate the signature surface, which needs the types only, and they miss generated constants
without a source range. The modules under \lean{Init} and \lean{Std} form a library of basic data
structures, and a library of mathematics may behave differently. The sample is systematic in the
order of names, so it is spread over the namespaces without being random.

%% 6 ---------------------------------------------------------------------------------------
\section{What the Calculus Leaves Out}
\label{sec:limits}

The calculus is untyped and equational. A proof assistant differs in ways that bear on each
result.

\emph{Types.} In a dependent type theory a constant has a type, and a declared constant is an
axiom by itself, since it asserts that its type is inhabited. The audit surface then holds the
type of every constant it reaches, as in Section~\ref{sec:measure}. Theorem~\ref{thm:audit}
needs weakening, substitution of definitions, and cut for the type theory in question. We expect
these to hold for the theories whose kernels have been formalized~\citep{sozeau2020,carneiro2024}
and have not proved it.

\emph{Elaboration.} The audit surface consists of kernel terms, and a human reads source text.
Notations, implicit arguments, coercions, and instances decide which kernel term a line of
source denotes, and a system may print a term in a way that reads as a different
statement~\citep{pollack1998,wiedijk2012}. The audit surface says which declarations matter. It
does not make them easy to read.

\emph{Kernel extensions.} A kernel may evaluate some functions with built-in code and may accept
the result of compiled code as an axiom. Such an axiom appears in the audit surface, and the code
it trusts lies outside the library.

\emph{Global conditions.} Theorem~\ref{thm:merge} relies on well-formedness conditions that
concern one entry and the entries before it. A kernel that keeps one graph of universe constraints
for all declarations checks a global condition, and two libraries can each satisfy it while their
union does not. Instances and notations do not affect checked terms, and they can change what
the source of a later entry means.

\emph{The kernel.} The audit surface assumes a correct kernel. Verified kernels address this
assumption~\citep{harrison2006,kumar2014,candle2022,sozeau2020,carneiro2024}.

%% 7 ---------------------------------------------------------------------------------------
\section{Related Work}
\label{sec:related}

\emph{Believing a checked proof.} \citet{lcf1979} introduced the architecture in which only
inference rules construct theorems, and the de Bruijn criterion asks for proofs that a small
independent program can check~\citep{barendregt2005,geuvers2009}. \citet{pollack1998} separates
two questions about a machine-checked proof, whether it is a proof and what it proves, and
\citet{wiedijk2012} studies systems in which the printed form of a statement misleads.
\citet{adams2016} describes how to audit a formalization such as the proof of the Kepler
conjecture~\citep{hales2017}: an auditor reviews the statement and the definitions it depends on
and replays the proof in a trusted checker. Our audit surface makes the reviewed set precise for
libraries with axioms and proves that it suffices.

\emph{Small trusted bases.} Proof-carrying code checks the proofs of an untrusted
producer~\citep{necula1997}, and its foundational form reduces what the checker
trusts~\citep{appel2001,appel2003}. Verified kernels reduce the trust in the checker
itself~\citep{harrison2006,kumar2014,candle2022,sozeau2020,carneiro2024}. These works concern the
kernel. We take a correct kernel as given and ask what remains for the reader.

\emph{Definitions.} \citet{kumar2014} and \citet{arthan2014} study the principles by which HOL
introduces constants and prove them sound. \citet{severi1994} add definitions to pure type
systems. \citet{gratzer2022} give a type theory in which the unfolding of a definition is
controlled. Definition elimination in Theorem~\ref{thm:audit} is the classical conservativity of
definitions, applied selectively: the definitions in the audit set stay because the statement or
an axiom depends on them.

\emph{Modules and linking.} Our libraries are signatures in the sense of logical
frameworks~\citep{harper1993}, extended with definitions and proofs. Theorem~\ref{thm:merge}
corresponds to linking separately checked fragments~\citep{cardelli1997}. The distinction between
transparent and sealed definitions in Theorem~\ref{thm:replace} is that of manifest and abstract
components of a module~\citep{leroy1994,mitchell1988}.

%% 8 ---------------------------------------------------------------------------------------
\section{Conclusion}

A machine-checked theorem asks a human to read its audit surface: the definitions under the
statement and the axioms under the proof. The statement follows from the surface alone, and the
surface is found in one pass. In Lean's library the surface is about a third of what the kernel
checked, and it is about three orders of magnitude more than the statement.

The results are proved for an untyped equational calculus. The step to a dependent type theory
with inductive types is open, and so is a surface defined on source text. The measurement
suggests two further directions. A reader needs a way to name the subjects of a theorem so that
the surface excludes their bodies soundly. A library needs a measure of its total surface, which
the union in Section~\ref{sec:measure} approximates.

\bibliographystyle{ACM-Reference-Format}
\bibliography{refs}

\end{document}
